%_________Appendix__________________________________

% final tutorial remarks
\ifLSRITRtutorial
	\chapter{BUSTED! Last chance to actually read the HowTo-Section!}
	
%	Make sure your thesis is well structured, that each major section does what it is supposed to do, and that the whole thing hangs together. The basic structure is often as given in this template (but other structures are possible). In particular, don't think you need to have exactly as many major sections or chapters as the list implies; sometimes it makes sense to merge things, sometimes it makes sense to move things (e.g., the literature review is in many papers deferred until after the results), sometimes it makes sense to split a logical part into several individual sections. Just use some common sense.
%	
%	Hand in your thesis at minimum \textbf{one week} before the deadline for correction. You will receive feedback for the final version and very likely have to do minor or major revisions of your writing. Plan your writing schedule to allow for these adjustments, which can have quite some impact on your grade! 
%	
%	\optional{Please have a look on our \href{https://wiki.lsr.ei.tum.de/thesiswriting_students}{thesis-guidelines} as well before submitting your \emph{final} thesis.}
	
\optional{Do not forget to check our \href{https://wiki.tum.de/display/lsritr/Students}{\underline{thesis-guidelines}} before submitting your \emph{final} thesis.}	
	
%	\section{Style and Expressions}
%	
%	Before handing in your thesis, even for an intermediate review, please perform a spellcheck and correct grammar mistakes. The report is not meant to be a narrative text. Please stick to neutral and technical style and avoid subjective or biased expressions or adjectives/adverbs such as \emph{obviously, always, very, especially well, actually, so-called etc}. Scientific writing is about precision and you should underpin your statements factually, not soften them with unnecessary qualifiers.
%	
%	... Okay enough. But please check chapter \ref{sec:Tutorial} before starting with your report. 

\section{Spellcheck reminder}
Before handing in your thesis, even for an intermediate review, please perform a proper spellcheck and correct grammar mistakes. Otherwise, there is a high likelihood that your supervisor will return your report immediately without making any comments.
\fi

\chapter{Experimental Details}
\label{appendix:experimental_details}

\section{Acrobot Evaluation Models}
\label{appendix:acrobot_evaluation_models}

\subsection{Reduced Maximal-Coordinates LGVI Model}
\label{appendix:reduced_maximal_coordinates_LGVI_model}

\textbf{Approximation used.} This model is algebraically equivalent to Model 1 in the exact polynomial problem, but uses the kinematics to remove \(R_{k-1}\) from the dynamics clique. The neighboring rotations are replaced by
\[
R_{i,k+1}=R_{i,k}F_{i,k},
\qquad
R_{i,k-1}=R_{i,k}F_{i,k-1}^{\top}.
\]
Hence
\[
\widehat\Delta^s_{i,k}
=s_{i,k}(a_{i,k}+a_{i,k-1}-2)+c_{i,k}(b_{i,k}-b_{i,k-1}),
\]
\[
\widehat\Delta^c_{i,k}
=c_{i,k}(a_{i,k}+a_{i,k-1}-2)-s_{i,k}(b_{i,k}-b_{i,k-1}).
\]
These are exactly the same second differences as in Model 1, only rewritten with \(R_k,F_{k-1},F_k\).

\textbf{Reduced translational dynamics.}
\[
m_1\frac{\ell_1}{2}\widehat\Delta^s_{1,k}
-h^2\Big(\bar\lambda^{(1)}_{0,k}+\bar\lambda^{(1)}_{12,k}\Big)=0,
\]
\[
-m_1\frac{\ell_1}{2}\widehat\Delta^c_{1,k}
+h^2m_1g
-h^2\Big(\bar\lambda^{(2)}_{0,k}+\bar\lambda^{(2)}_{12,k}\Big)=0,
\]
\[
m_2\left[\ell_1\widehat\Delta^s_{1,k}+\frac{\ell_2}{2}\widehat\Delta^s_{2,k}\right]
+h^2\bar\lambda^{(1)}_{12,k}=0,
\]
\[
-m_2\left[\ell_1\widehat\Delta^c_{1,k}+\frac{\ell_2}{2}\widehat\Delta^c_{2,k}\right]
+h^2m_2g+h^2\bar\lambda^{(2)}_{12,k}=0.
\]

\textbf{Clique choice.} Since \(R_{k-1}\) does not appear explicitly after substitution, one interior stage can be covered by
\[
\boxed{\mathcal C^{21}_k=
\{R_k,R_{k+1},F_{k-1},F_k,\lambda_{0,k},\lambda_{12,k},u_k\}},
\qquad k=1,\dots,N-1.
\]
The size is
\[
2\cdot4+2\cdot4+4+1=\boxed{21}.
\]
Boundary and terminal cliques remain
\[
\mathcal C_0=\{R_0,R_1, R_2, F_0, F_1, \lambda_{0,1},\lambda_{12,1},u_1\},\qquad |\mathcal C_0|=25,
\]
\[
\mathcal C_N=\{R_{N-1}, R_N, F_{N-2}, F_{N-1}, \lambda_{0,N-1},\lambda_{12,N-1},u_{N-1}\},\qquad |\mathcal C_N|=21.
\]
The total number of variables is again \(13N-1\). The advantage is the smaller clique. The drawback is that the sparse relaxation can become weaker,
because \(R_{k-1}\) is not in the same moment block as the rest of the three-node stencil.

\subsection{Discretized Newton-Euler-Lagrange Model}
\label{appendix:discretized_newton_euler_lagrange_model}

\paragraph{Coordinates and approximation.}
Let
\[
q_k=
\begin{bmatrix}q_{1,k}\\q_{2,k}\end{bmatrix},
\qquad
\dot q_k=\omega_k=
\begin{bmatrix}\omega_{1,k}\\\omega_{2,k}\end{bmatrix},
\]
where $q_{1,k}$ is the absolute shoulder angle and $q_{2,k}$ is the
relative elbow angle. Thus, $R_{2,k}$ represents the relative joint
coordinate, not the absolute orientation of the second link. For
$i\in\{1,2\}$, define
\[
R_{i,k}=
\begin{bmatrix}
c_{i,k}&-s_{i,k}\\
s_{i,k}&c_{i,k}
\end{bmatrix},
\qquad
F_{i,k}=
\begin{bmatrix}
a_{i,k}&-b_{i,k}\\
b_{i,k}&a_{i,k}
\end{bmatrix}.
\]
Over the interval $[t_{k-1},t_k]$, the torque and angular velocity are
evaluated at the left endpoint:
\[
u(t)\approx u_{k-1},\qquad
\dot q(t_{k-1})=\omega_{k-1},\qquad
\ddot q(t_{k-1})
\approx\frac{\omega_k-\omega_{k-1}}{h}.
\]
The Lie-group reconstruction is
\[
R_{i,k}=R_{i,k-1}F_{i,k-1},
\]
or equivalently
\[
c_{i,k}-c_{i,k-1}a_{i,k-1}+s_{i,k-1}b_{i,k-1}=0,
\]
\[
s_{i,k}-s_{i,k-1}a_{i,k-1}-c_{i,k-1}b_{i,k-1}=0.
\]
Following the cubic approximation, the step sine is related
to the physical angular velocity by
\[
b_{i,k-1}
-h\omega_{i,k-1}
+\frac{h^3}{6}\omega_{i,k-1}^3=0.
\]
The associated polynomial rotation constraints are
\[
c_{i,k}^2+s_{i,k}^2=1,\qquad
a_{i,k}^2+b_{i,k}^2=1,\qquad
a_{i,k}\geq a_{\min}>0.
\]

\paragraph{Discrete Newton--Euler--Lagrange dynamics.}
The mass matrix, Coriolis vector, and gravity torque evaluated at node
$k$ are
\[
M_k=
\begin{bmatrix}
I_1+I_2+m_2l_1^2+2m_2l_1l_{c2}c_{2,k}
&
I_2+m_2l_1l_{c2}c_{2,k}
\\
I_2+m_2l_1l_{c2}c_{2,k}
&
I_2
\end{bmatrix},
\]
\[
\gamma_k:=C(q_k,\omega_k)\omega_k
=
\begin{bmatrix}
-m_2l_1l_{c2}s_{2,k}
\left(
2\omega_{1,k}\omega_{2,k}
+\omega_{2,k}^2
\right)
\\
m_2l_1l_{c2}s_{2,k}\omega_{1,k}^2
\end{bmatrix},
\]
\[
\tau_{g,k}
=
\begin{bmatrix}
-m_1gl_{c1}s_{1,k}
-m_2g\!\left[
l_1s_{1,k}
+l_{c2}
\left(
s_{1,k}c_{2,k}
+c_{1,k}s_{2,k}
\right)
\right]
\\
-m_2gl_{c2}
\left(
s_{1,k}c_{2,k}
+c_{1,k}s_{2,k}
\right)
\end{bmatrix},
\qquad
B=
\begin{bmatrix}
0\\1
\end{bmatrix}.
\]
Hence, for $k=1,\ldots,N$, the left-endpoint polynomial dynamics are
\[
\boxed{
M_{k-1}\left(\omega_k-\omega_{k-1}\right)
+hC(q_{k-1},\omega_{k-1})\omega_{k-1}
-h\tau_{g,k-1}
-hBu_{k-1}=0.
}
\]
These two scalar equations, together with the reconstruction, cubic-step,
and circle constraints, form the discrete transition model. The highest
polynomial degree is three.

\paragraph{Trajectory optimization and chain clique.}
Define the augmented polynomial state
\[
x_k=
\begin{bmatrix}
c_{1,k}&s_{1,k}&a_{1,k}&b_{1,k}&\omega_{1,k}&
c_{2,k}&s_{2,k}&a_{2,k}&b_{2,k}&\omega_{2,k}
\end{bmatrix}^{\!T}\in\mathbb{R}^{10}.
\]

The standard interior chain clique is
\[
\boxed{
z(I_k)=\{x_{k-1},u_{k-1},x_k\},
\qquad |I_k|=10+1+10=21.
}
\]

\section{Model Evaluation}
\label{appenidx:model_evaluation}

\subsection{Numerical Simulation Validation}
\label{appendix:numerical_simulation_validation}

\subsection{SDP Tightness and Dynamic Residuals}
\label{appendix:SDP_tightness_and_dynamic_residuals}
